\chapter{Quotients and Algebraic Hierarchy}

% Chapter orientation paragraph
In this chapter, we address a fundamental limitation of the
\lstinline[language=lean]|myPreRat| structure from the previous
chapter: distinct representations of the same rational number are
not identified as equal.
We resolve this using quotient types
(Section~\ref{sec:quotients}), then step back to examine how
Mathlib's algebraic hierarchy is built using the type class
mechanism introduced earlier (Section~\ref{sec:typeclasses}).

\section{Quotients}\label{sec:quotients}

% FIXED: "gopes" → "goes"
Let us look at what goes wrong with \lstinline[language=lean]|myPreRat|:
\begin{lstlisting}[language=lean]
example : myPreRat.mk 2 4 (by decide) ≠ myPreRat.mk 1 2 (by decide) := by
  simp
\end{lstlisting}
This example proves that \lstinline[language=lean]|myPreRat.mk 2 4|
and \lstinline[language=lean]|myPreRat.mk 1 2| are not equal, even
though mathematically $\frac{2}{4} = \frac{1}{2}$!
The name \lstinline[language=lean]|myPreRat| was already alluding to
the need for further work.

In mathematics, we treat different representations of the same
rational number as equivalent through an
\textbf{equivalence relation}~(\cite{algebrology_rationals_2019}).
We consider fractions like $\frac{1}{2}$, $\frac{2}{4}$, and
$\frac{3}{6}$ as belonging to the same equivalence class.
More precisely, mathematics uses \textbf{quotients} to group elements
of a set by an equivalence relation.
For instance, the equivalence class
$\left[\frac{1}{2}\right] = \left\{\ldots, \frac{-2}{-4},
  \frac{1}{2}, \frac{2}{4}, \frac{3}{6}, \ldots\right\}$
represents all fractions equivalent to $\frac{1}{2}$.
The set of rational numbers is the set of such equivalence classes.

% FIXED: P33 — Z* means invertible elements {1,-1}, not Z \ {0}.
Algebraically, each rational number can be represented as a pair of
integers $(a, b)$ where $b \neq 0$.
This construction is formalized by an equivalence relation
$\sim_{\mathbb{Q}}$:


$$
  \mathbb{Q} = \bigl(\mathbb{Z} \times
    (\mathbb{Z} \setminus \{0\})\bigr) \big/ {\sim_{\mathbb{Q}}}
$$


The equivalence relation is defined as:


$$
  (a,b) \sim_{\mathbb{Q}} (c,d) \iff ad = bc \quad
  \text{for all } (a,b),\, (c,d) \in
  \mathbb{Z} \times (\mathbb{Z} \setminus \{0\})
$$



The relation $\sim_{\mathbb{Q}}$ must satisfy reflexivity, symmetry,
and transitivity.
Moreover, each operation defined on the quotient set must be
\textbf{well-defined}; that is, the result must not depend on the
choice of representatives.
If we naively define addition as


$$
  (a,b) + (c,d) = (a+c,\, b+d),
$$


we encounter a problem.
Consider:


$$
  (1,2) + (1,3) = (2,5) \quad \text{and} \quad
  (1,2) + (2,6) = (3,8).
$$


But $(1,3) \sim_{\mathbb{Q}} (2,6)$ since $1 \cdot 6 = 2 \cdot 3$,
yet $(2,5) \not\sim_{\mathbb{Q}} (3,8)$ since
$2 \cdot 8 = 16 \neq 15 = 5 \cdot 3$.
A well-defined definition for addition is instead:


$$
  (a,b) + (c,d) = (ad+bc,\, bd).
$$


We will verify this in our formalization.

Similarly, in type theory and Lean we have \textbf{quotient types},
which allow us to construct new types by means of equivalence
relations.
We now define proper rational numbers using quotient types in Lean.
As mentioned earlier, this differs from Mathlib's approach, which
achieves the same goal by mechanically reducing each rational number
to a canonical form using coprimality.

We start from the notion of equivalence relations.
The structure \lstinline[language=lean]|Equivalence| has fields
\lstinline[language=lean]|refl|, \lstinline[language=lean]|symm|,
and \lstinline[language=lean]|trans|.
You can inspect its definition using
\lstinline[language=lean]|#print Equivalence|.
The \lstinline[language=lean]|Setoid| type class encapsulates an
equivalence relation on a given type, requiring that the relation is
indeed an equivalence via its \lstinline[language=lean]|iseqv| field.

% REVISED: P35 — present well-definedness proofs BEFORE the instances
% that use them, so dependencies are satisfied in reading order.
% We first define the Setoid, then the well-definedness proofs,
% then the quotient operations.

We begin by defining the \lstinline[language=lean]|Setoid| instance
and the quotient type:

\begin{lstlisting}[language=lean]
instance myRel : Setoid myPreRat where
  r p q := p.num * q.den = q.num * p.den
  iseqv := by
    constructor
    · intro p; rfl
    · rintro ⟨p, p', hp'⟩ ⟨q, q', hq'⟩
      simp [Eq.comm]
    · rintro ⟨p, p', hp'⟩ ⟨q, q', hq'⟩ ⟨r, r', hr'⟩ hpq hqr
      apply mul_left_cancel₀ (mod_cast hq'.ne' : q' ≠ (0 : ℤ))
      grind

abbrev myRat : Type := Quotient myRel
\end{lstlisting}

% FIXED: P34 — removed unnecessary simp_all line
In the \lstinline[language=lean]|iseqv| field, we prove reflexivity,
symmetry, and transitivity of our relation.
Reflexivity (\lstinline[language=lean]|intro p; rfl|) is immediate.
Symmetry uses \lstinline[language=lean]|Eq.comm| to swap the sides of
the equation.
For transitivity, the key step involves canceling the common factor
\lstinline[language=lean]|q.den| using
\lstinline[language=lean]|mul_left_cancel₀| (left multiplication
cancellation for integers), which states that if
$a \cdot b = a \cdot c$ and $a \neq 0$, then $b = c$.
The expression
\lstinline[language=lean]|mod_cast hq'.ne' : q' ≠ (0 : ℤ)| takes
\lstinline[language=lean]|hq'| (which states
\lstinline[language=lean]|0 < q.den| as a natural number) and casts
it to a proof that \lstinline[language=lean]|q.den ≠ 0| as an
integer.
Finally, \lstinline[language=lean]|abbrev| is a Lean keyword that
creates a type abbreviation.
Unlike \lstinline[language=lean]|def|, which creates a new definition
that must be unfolded explicitly, with
\lstinline[language=lean]|abbrev| Lean automatically treats
\lstinline[language=lean]|myRat| as
\lstinline[language=lean]|Quotient myRel| without requiring manual
unfolding.

Before defining operations on \lstinline[language=lean]|myRat|, we
must prove that our operations are well-defined on the quotient, that
is, that the result does not depend on the choice of representative.
This is achieved using \lstinline[language=lean]|Quotient.lift₂|: to
define a function from a quotient type, we must provide an underlying
function on the representatives and a proof that equivalent
representatives yield the same result.

We first prove that the ordering respects the equivalence relation.
The heart of the proof is the forward direction:

\begin{lstlisting}[language=lean]
private theorem le_respects_equiv_forward
    (a₁ b₁ a₂ b₂ : myPreRat)
    (ha : a₁ ≈ a₂) (hb : b₁ ≈ b₂)
    (h : a₁ ≤ b₁) : a₂ ≤ b₂ := by
  have pos_prod : (0 : Int) < (a₁.den * b₁.den) :=
    myPreRat.den_prod_pos a₁ b₁
  have pos_prod2 : 0 < (a₂.den * b₂.den : Int) :=
    myPreRat.den_prod_pos a₂ b₂
  apply Int.le_of_mul_le_mul_right _ pos_prod
  calc (a₂.num * b₂.den) * (a₁.den * b₁.den)
      = a₂.num * a₁.den * b₂.den * b₁.den := by ring
    _ = a₁.num * a₂.den * b₂.den * b₁.den := by rw [← ha]
    _ = a₁.num * b₁.den * (a₂.den * b₂.den) := by ring
    _ ≤ b₁.num * a₁.den * (a₂.den * b₂.den) :=
        Int.mul_le_mul_of_nonneg_right h (Int.le_of_lt pos_prod2)
    _ = b₁.num * b₂.den * a₁.den * a₂.den := by ring
    _ = b₂.num * b₁.den * a₁.den * a₂.den := by rw [← hb]
    _ = (b₂.num * a₂.den) * (a₁.den * b₁.den) := by ring
\end{lstlisting}

Given $h : a_1 \leq b_1$ (meaning
$a_1.\text{num} \cdot b_1.\text{den} \leq
  b_1.\text{num} \cdot a_1.\text{den}$) and equivalences
$ha : a_1 \approx a_2$ and $hb : b_1 \approx b_2$, we need to prove
$a_2 \leq b_2$.
The strategy is to introduce a common positive factor and use the
given information.
We apply \lstinline[language=lean]|Int.le_of_mul_le_mul_right|, which
states that to prove $X \leq Y$, it suffices to prove
$X \cdot Z \leq Y \cdot Z$ for positive $Z$, then cancel $Z$.
We choose $Z = a_1.\text{den} \cdot b_1.\text{den}$ (shown positive
by \lstinline[language=lean]|pos_prod|).
The \lstinline[language=lean]|calc| block proceeds as follows: we
rearrange the left side and substitute using $ha$, apply $h$
(multiplying by the positive factor $a_2.\text{den} \cdot
  b_2.\text{den}$ to preserve the inequality), rearrange and
substitute using $hb$, and finally rearrange to match the required
form.
After canceling the common factor, we obtain $a_2 \leq b_2$.

Throughout the proof, we use the \lstinline[language=lean]|ring|
tactic from Mathlib, which automates proofs of equalities in
commutative rings (such as $\mathbb{Z}$).

The full well-definedness proof for ordering uses symmetry to avoid
duplication:

\begin{lstlisting}[language=lean]
theorem myRel_respects_le (a₁ b₁ a₂ b₂ : myPreRat) :
    a₁ ≈ a₂ → b₁ ≈ b₂ → (a₁ ≤ b₁) = (a₂ ≤ b₂) := by
  intro ha hb
  simp only [eq_iff_iff]
  constructor
  · exact le_respects_equiv_forward a₁ b₁ a₂ b₂ ha hb
  · exact fun h =>
      le_respects_equiv_forward a₂ b₂ a₁ b₁ ha.symm hb.symm h
\end{lstlisting}

We transform the equality of propositions into a biconditional using
\lstinline[language=lean]|eq_iff_iff|, then prove both directions
with \lstinline[language=lean]|constructor|.
Since both goals are symmetric, we reuse
\lstinline[language=lean]|le_respects_equiv_forward|.

We also prove that addition is well-defined:

\begin{lstlisting}[language=lean]
theorem myRel_respects_add (a₁ b₁ a₂ b₂ : myPreRat) :
    a₁ ≈ a₂ → b₁ ≈ b₂ → (a₁ + b₁) ≈ (a₂ + b₂) := by
  intro ha hb
  calc (a₁.num * b₁.den + b₁.num * a₁.den) * (a₂.den * b₂.den)
      = a₁.num * b₁.den * a₂.den * b₂.den
        + b₁.num * a₁.den * a₂.den * b₂.den := by ring
    _ = a₂.num * a₁.den * b₁.den * b₂.den
        + b₂.num * b₁.den * a₁.den * a₂.den := by rw [← ha, ← hb]; ring
    _ = (a₂.num * b₂.den + b₂.num * a₂.den) * (a₁.den * b₁.den)
        := by ring
\end{lstlisting}

We need to show that if $a_1 \approx a_2$ and $b_1 \approx b_2$,
then $(a_1 + b_1) \approx (a_2 + b_2)$.
By unfolding the addition definition and the relation, we must prove:
\begin{align*}
       & (a_1.\text{num} \cdot b_1.\text{den} +
  b_1.\text{num} \cdot a_1.\text{den}) \cdot
  (a_2.\text{den} \cdot b_2.\text{den})         \\
  = \; & (a_2.\text{num} \cdot b_2.\text{den} +
  b_2.\text{num} \cdot a_2.\text{den}) \cdot
  (a_1.\text{den} \cdot b_1.\text{den})
\end{align*}
The \lstinline[language=lean]|calc| proof proceeds in three steps.
First, we distribute the product over the sum using
\lstinline[language=lean]|ring|.
Next, we apply the equivalences $ha$ and $hb$ using
\lstinline[language=lean]|rw [← ha, ← hb]|, which substitutes
$a_1.\text{num} \cdot a_2.\text{den}$ with
$a_2.\text{num} \cdot a_1.\text{den}$ and similarly for $b$.
Finally, we factor back into sum-times-product form using
\lstinline[language=lean]|ring|.

Now that we have established well-definedness, we can define
operations on the quotient type \lstinline[language=lean]|myRat|:
\newpage
\begin{lstlisting}[language=lean]
instance : LE myRat where
  le r₁ r₂ := Quotient.lift₂ (fun a b ↦ a ≤ b)
    myRel_respects_le r₁ r₂

instance : Add myRat where
  add r₁ r₂ := Quotient.lift₂ (fun a b ↦ ⟦a + b⟧)
    (fun a₁ b₁ a₂ b₂ ha hb ↦
      Quotient.sound (myRel_respects_add a₁ b₁ a₂ b₂ ha hb))
    r₁ r₂

instance : OfNat myRat 0 where
  ofNat := ⟦myPreRat.zero⟧

lemma add_nonneg (a b : myRat) :
    0 ≤ a → 0 ≤ b → 0 ≤ a + b := by
  induction a using Quotient.ind with | _ a =>
  induction b using Quotient.ind with | _ b =>
  intro ha hb
  exact myPreRat.add_nonneg a b ha hb
\end{lstlisting}

Notice how we reuse the previous
\lstinline[language=lean]|add_nonneg| lemma for the quotient type.
The syntax
\lstinline[language=lean]|induction a using Quotient.ind with _ a =>|
unwraps the quotient value \lstinline[language=lean]|a : myRat| to
its underlying representative
\lstinline[language=lean]|a : myPreRat|.

We have now constructed
\lstinline[language=lean]|myRat.add_nonneg|.
However, our original discussion was about proving transitivity of
the less-or-equal operator.
In our earlier work with natural numbers, we used
\lstinline[language=lean]|Nat.le_trans|, a theorem specifically for
natural numbers.
However, the transitivity property holds not only for naturals but
also for integers, reals, and any partially ordered set.

% REVISED: P38 — corrected explanation: le_trans is not doing a clever
% general argument. It's about consistent interface, not avoiding
% duplication in *proving* transitivity.
Rather than having each piece of downstream code call a different
type-specific lemma (\lstinline[language=lean]|Nat.le_trans|,
\lstinline[language=lean]|Rat.le_trans|, etc.), Mathlib provides a
general lemma \lstinline[language=lean]|le_trans| that works for any
type endowed with a partial order.
It is important to understand what this generality achieves: the
general \lstinline[language=lean]|le_trans| is not performing a
clever argument---it is simply retrieving the transitivity proof that
was supplied as part of the partial order structure.
We still need to provide a new proof of transitivity whenever we
create a new \lstinline[language=lean]|PartialOrder| instance.
The benefit is \emph{downstream}: any code that uses transitivity
can call the single lemma \lstinline[language=lean]|le_trans| rather
than needing to know which specific type it is working with.
This provides a consistent interface for all ordered types.

Mathlib achieves this through type classes and a carefully
constructed algebraic hierarchy.
We already used type classes such as \lstinline[language=lean]|Add|
and \lstinline[language=lean]|LE| to define operations on
\lstinline[language=lean]|myPreRat|.
Knowing that rational numbers form a totally ordered set, we can
enhance \lstinline[language=lean]|myRat| with the
\lstinline[language=lean]|LinearOrder| type class:

% ADDED: P40 — note about Mathlib version compatibility
\begin{remark}
  The following code was developed and tested with Mathlib version
  4.24.0.
  Later versions of Mathlib may introduce instances (such as
  \lstinline[language=lean]|Quotient.instPreorder|) that conflict
  with the order structure defined here.
  This is a consequence of Mathlib's ongoing development and not a
  fundamental issue with the construction.
\end{remark}

\begin{lstlisting}[language=lean]
noncomputable instance : LinearOrder myRat where
  le_refl p := by
    induction p using Quotient.ind with | _ a =>
    exact myPreRat.le_refl a

  le_trans p q r := by
    induction p using Quotient.ind with | _ a =>
    induction q using Quotient.ind with | _ b =>
    induction r using Quotient.ind with | _ c =>
    intro hab hbc
    exact myPreRat.le_trans a b c hab hbc

  le_antisymm p q := by
    induction p using Quotient.ind with | _ a =>
    induction q using Quotient.ind with | _ b =>
    intro hab hba
    exact Quotient.sound (myPreRat.le_antisymm a b hab hba)

  le_total p q := by
    induction p using Quotient.ind with | _ a =>
    induction q using Quotient.ind with | _ b =>
    cases Int.le_total (a.num * b.den) (b.num * a.den) with
    | inl h => exact Or.inl h
    | inr h => exact Or.inr h

  toDecidableLE := Classical.decRel _
end myRat
\end{lstlisting}

The structure follows the same pattern: use
\lstinline[language=lean]|Quotient.ind| to unwrap quotient values to
their representatives, then apply the corresponding proof for
\lstinline[language=lean]|myPreRat|.
The antisymmetry case uses
\lstinline[language=lean]|Quotient.sound|, which states that if two
representatives are equivalent ($a \approx b$), then their
equivalence classes are equal
($\llbracket a \rrbracket = \llbracket b \rrbracket$).

\newpage
The underlying proofs for \lstinline[language=lean]|myPreRat| are:
\begin{lstlisting}[language=lean]
theorem le_refl (a : myPreRat) : a ≤ a := by
  exact Int.le_refl _

theorem le_trans (a b c : myPreRat) :
    a ≤ b → b ≤ c → a ≤ c := by
  intro hab hbc
  apply Int.le_of_mul_le_mul_right _ b.den_pos_int
  calc (a.num * c.den) * b.den
      = (a.num * b.den) * c.den := by ring
    _ ≤ (b.num * a.den) * c.den :=
        Int.mul_le_mul_of_nonneg_right hab
          (Int.le_of_lt c.den_pos_int)
    _ = (b.num * c.den) * a.den := by ring
    _ ≤ (c.num * b.den) * a.den :=
        Int.mul_le_mul_of_nonneg_right hbc
          (Int.le_of_lt a.den_pos_int)
    _ = (c.num * a.den) * b.den := by ring

theorem le_antisymm (a b : myPreRat) :
    a ≤ b → b ≤ a → a ≈ b := by
  intro hab hba
  exact Int.le_antisymm hab hba
\end{lstlisting}

Reflexivity and antisymmetry are straightforward applications of the
corresponding integer properties.
For transitivity, given $hab : a \leq b$ and $hbc : b \leq c$, we
need to prove $a \leq c$.
The strategy is to introduce a common positive factor
$b.\text{den}$, prove the inequality with this factor, then cancel it
using \lstinline[language=lean]|Int.le_of_mul_le_mul_right|.
The \lstinline[language=lean]|calc| chain rearranges terms, applies
$hab$ and $hbc$ (multiplying by positive factors to preserve the
inequalities), and rearranges to match the required form.

The property \lstinline[language=lean]|le_total| asserts that the
order is total: for any $a, b$, either $a \le b$ or $b \le a$.
By unwrapping the quotient, this reduces to an inequality of
integers, and we apply \lstinline[language=lean]|Int.le_total|.

The field
\lstinline[language=lean]|toDecidableLE := Classical.decRel _|
requires special attention.
In Mathlib, a \lstinline[language=lean]|LinearOrder| is not merely a
set of axioms; it also defines
\lstinline[language=lean]|min| and \lstinline[language=lean]|max|
functions.
By default, \lstinline[language=lean]|min a b| is defined as
\lstinline[language=lean]|if a ≤ b then a else b|.
To evaluate this conditional, Lean requires the relation $\le$ to be
\emph{decidable} (i.e., there must be an algorithm to determine
whether it holds).
While we could write an explicit comparison algorithm for our
rationals, we opted for a concise approach using
\lstinline[language=lean]|Classical.decRel|, which invokes classical
logic to treat every proposition as decidable.
Because this relies on the axiom of choice rather than an executable
algorithm, the resulting functions cannot be computed by the kernel,
necessitating the \lstinline[language=lean]|noncomputable| marker.

Here is the full construction of \lstinline[language=lean]|myRat|:
[\href{https://live.lean-lang.org/#codez=JYWwDg9gTgLgBAWQIYwBYBtgCMBQODOMUArgMYzFQCmcIAngArUBKKcA7qldTnHAHbEQcAFxwAkvxi84AEyr9RcAHIoZ8/gH1I+JQAY4AHjkK8/JCCr4wSUjXpMqraTmD9CSfnaUAZAKK0jCxsnNxUMug0UICBBHBQgEEEogC8cdEAdILCAFRwgImECWkacIAmRHHxGUJwOfnpGnhuHl40YgCCsrKBjs4cXDx8SO2pZclwAN4yfJkjMRXZeQVFANRls1XztQoANBMmiiIpM0U5C1s7GtoQuvsqKGkgxOgXuocKTyvnOjIAvnjyAGZwABe3AgSgcwXg11GAkq1z0m12IwAjAiPpcRlg6CZSMB5HAfg0YJ5vGIAPJ/VTwcFONgGUJ9OAQClsa7AqAQPCREAgJACCD8fhUADmmmAfwBAAooGCgjSYABKfQlOJwQAphHADKUoKtrpiZKR+YQSORoEYANwAPjgbiIoNQMgA7XBSKhPEKaAZjoUFMrtVMcki4GxUGa4PhQGAg/AQ3AqAAPWzRx3O138d0atbaopa1YB0Ph8Ch+OJuD2nBcnlB9qafj8wVCuAS3lYGVdFCKsSaoNwQBJhBnSi2+13ecsW7q6DIC5H+egsQBtWsC4Wi8UAXRkNvZpd5qFw/TAYFnEikaQGshrdeFyeL5GPMDuDwvS4brsbknv5hgAGEkIQn/W4Cwb1+Hla8E1vd8H0eRcAN3N8T0/H8/xg4Ug2A0CgA}{link to Lean live}]

% ADDED: transition
With the use of type classes such as
\lstinline[language=lean]|LinearOrder|, we enhanced
\lstinline[language=lean]|myRat| so that general theorems about
linear orders apply to it automatically.
In the next section, we examine this mechanism more carefully and
discuss Mathlib's broader algebraic hierarchy.

\section{Type Classes and Algebraic Hierarchy}\label{sec:typeclasses}

Structures and type classes are heavily used in Mathlib for providing
properties such as ordered rings, fields, metric spaces, and
topological spaces, which we refer to as
\textbf{algebraic structures}.
An algebraic structure is defined by a carrier type and a set of
operations and properties.
We write $C\;\alpha$ where $\alpha$ is the \textbf{carrier type}
(the underlying set of elements) and $C$ is the \textbf{structure}
(the operations and properties defined on that type).
For example, in a group structure
$\text{Group}\;\alpha$, $\alpha$ is the carrier type (such as
$\mathbb{Z}$ or $\mathbb{Q}$), and $\text{Group}$ defines the
multiplication, identity element, and inverse operation on $\alpha$.

% REVISED: P40 — noted that type classes can take multiple type arguments
Note that this is a simplified picture: type classes can take multiple
type arguments or additional non-type arguments.
For instance, \lstinline[language=lean]|Module ℚ ℝ| takes two types
(\lstinline[language=lean]|ℚ| as the scalar ring and
\lstinline[language=lean]|ℝ| as the module), and
\lstinline[language=lean]|OfNat α n| takes a type and a natural
number.

Under the hood, a type class is a structure.
The main difference between the two is that type classes use
\textbf{instance resolution}: a mechanism that allows Lean to
automatically infer properties about a type when the latter is
declared as an instance of a type class.
Using structures alone would require explicitly passing the structure
as an argument to functions that need it, which is cumbersome.

\newpage
Let us examine a simple example to see the main difference.
Using a structure:
\begin{lstlisting}[language=lean]
structure Semigroup' (α : Type*) where
  mul : α → α → α
  mul_assoc : ∀ a b c : α,
    mul (mul a b) c = mul a (mul b c)

def double {α : Type*} (s : Semigroup' α) (x : α) : α :=
  s.mul x x

def Semigroup'_Int : Semigroup' ℤ where
  mul := (· + ·)
  mul_assoc := Int.add_assoc

def Semigroup'_Rat : Semigroup' ℚ where
  mul := (· + ·)
  mul_assoc := Rat.add_assoc

#eval double Semigroup'_Int (-2)    -- -4
#eval double Semigroup'_Rat (1/2)   -- 1
\end{lstlisting}
The \lstinline[language=lean]|double| function takes an explicit
\lstinline[language=lean]|Semigroup' α| argument to perform the
doubling operation.
When calling \lstinline[language=lean]|double|, we must explicitly
provide the semigroup instance, such as
\lstinline[language=lean]|Semigroup'_Int| or
\lstinline[language=lean]|Semigroup'_Rat|.

In contrast, using a type class:
\begin{lstlisting}[language=lean]
class Semigroup'' (α : Type*) where
  mul : α → α → α
  mul_assoc : ∀ a b c : α,
    mul (mul a b) c = mul a (mul b c)

instance : Semigroup'' ℤ where
  mul := (· + ·)
  mul_assoc := Int.add_assoc

instance : Semigroup'' ℚ where
  mul := (· + ·)
  mul_assoc := Rat.add_assoc

def double' {α : Type*} [Semigroup'' α] (x : α) : α :=
  Semigroup''.mul x x

#eval double' (-2 : ℤ)    -- -4
#eval double' (1/2 : ℚ)   -- 1
\end{lstlisting}

Notice the square brackets
\lstinline[language=lean]|[Semigroup'' α]| in the definition of
\lstinline[language=lean]|double'|.
This syntax indicates that \lstinline[language=lean]|Semigroup'' α|
is a type class constraint.
When we call \lstinline[language=lean]|double'|, we do not need to
explicitly provide the semigroup instance.
Instead, Lean automatically infers the appropriate instance based on
the type of the argument.
This is the power of instance resolution.

We chose the semigroup example because it is simple yet demonstrates
how algebraic structures can be extended.
A semigroup can be extended to a monoid or a group by adding more
properties.
Using structures and type classes, we can implement this inheritance
pattern:

\begin{lstlisting}[language=lean]
class Group'' (α : Type*) extends Semigroup'' α where
  one : α
  left_id : ∀ a : α, mul one a = a
  right_id : ∀ a : α, mul a one = a
  inv : α → α
  left_inv : ∀ a : α, mul (inv a) a = one
  right_inv : ∀ a : α, mul a (inv a) = one

instance : Group'' ℤ where
  -- omitted

instance : Group'' ℚ where
  -- omitted

theorem mul_cancel₀ {α : Type*} [Group'' α] (a b c : α)
    (h : Semigroup''.mul a b = Semigroup''.mul a c) :
    b = c := by sorry
\end{lstlisting}

Now we can state the cancellation law once and use it for any type
that is an instance of \lstinline[language=lean]|Group''|.
We have already demonstrated this principle when discussing
transitivity of the less-or-equal relation and when enhancing
\lstinline[language=lean]|myRat| with the
\lstinline[language=lean]|LinearOrder| structure.

% ADDED: P42 — discussion of disadvantages of type class system
\paragraph{Limitations of type classes.}
While type classes are powerful, they are not without drawbacks.
Two important limitations are worth noting:
\begin{itemize}
  \item \textbf{Performance.}
        Instance resolution can be slow, especially in large projects
        like Mathlib where the hierarchy of type classes is deep and
        highly interconnected.
        Lean must search through many candidate instances to find the
        right one, and this search can become expensive.
  \item \textbf{Conflicting instances.}
        If two different instances of the same type class are registered
        for the same type, the results can be unpredictable or lead to
        subtle errors.
        For example, a type might carry two different orderings, but
        Lean's instance resolution will silently pick one, potentially
        leading to confusion.
        Mathlib takes great care to avoid such conflicts, but they can
        arise when combining independently developed libraries.
\end{itemize}

\paragraph{Mathlib's algebraic hierarchy.}
Mathlib builds an extensive algebraic hierarchy using type classes.
This hierarchy is far from being a simple tree; there are different
nuances and technical approaches that are beyond the scope of this
thesis.
For a detailed discussion on the design and implementation of
Mathlib's algebraic hierarchy, we refer the reader
to~\cite{mathlib2020}.

% ADDED: Chapter closing paragraph
In this chapter, we have seen how quotient types allow us to identify
equivalent representations, turning
\lstinline[language=lean]|myPreRat| into the well-behaved type
\lstinline[language=lean]|myRat|.
We then examined how Mathlib's algebraic hierarchy, built on type
classes and instance resolution, provides a consistent and extensible
framework for organizing mathematical structures.
In the next chapter, we turn to how topology and analysis are
formalized in Lean, building on the type class machinery introduced
here.